\chapter{World Model as Environment}
\label{chap:wm-env}

\section{Motivation}
\label{sec:wmep-motivation}

In regular offline-to-online reinforcement learning, the agent
is eventually allowed to interact with the real environment after being trained on a
fixed offline dataset. This second phase is used to correct the distributional
limitations of offline learning. Once online interaction begins, the policy can
collect fresh transitions, expand its coverage of the state space, and gradually
improve beyond what was observed in the dataset.

The main difficulty is that the
critic learned during offline training is only reliable on the support of the dataset.
As soon as the policy begins to drift away from that support during online fine-tuning,
the critic is forced to evaluate state-action pairs it has never seen. This leads to
uncontrolled value extrapolation, which the actor can then exploit. The resulting
feedback loop is well known: the policy moves into regions where the critic is
incorrect, the critic assigns high value to those regions due to function approximation
error, and the policy is further reinforced into them.

Existing methods address this instability by explicitly discouraging distributional
shift. Behaviour regularised approaches such as AWAC\cite{nair2021awacacceleratingonlinereinforcement} constrain
the policy to remain close to the dataset. Pessimistic methods such as Cal-QL\cite{nakamoto2024calqlcalibratedofflinerl} modify the critic targets to avoid overestimation outside
the data manifold. Model-based offline methods such as MOPO and COMBO\cite{yu2020mopomodelbasedofflinepolicy, yu2021combo} introduce explicit penalties
for uncertainty in learned dynamics. Although effective, these approaches share a
common principle: they attempt to keep learning anchored to the dataset rather than
fully replacing the environment during the online phase.

This chapter investigates a more direct alternative. Instead of collecting new
real-environment data during online training, we ask whether a latent
action-conditioned dynamics model can replace the environment entirely during this
phase. Concretely, we replace the simulator with the learned predictor
$f_\psi$ and define a set of substitution variants that test where such synthetic
interaction can break down. The variants are designed to separate representation
quality, predictor drift, value-target corruption, and uncertainty estimation before
the later chapters move to more constrained planning-based uses of the model.

% ---------------------------------------------------------------------
\section{Offline-to-Online Reinforcement Learning with World Model Encoder}
\label{sec:wmep-latent-obs}

Before replacing the environment with a world model, we first isolate the role of
representation learning. In particular, we ask whether the learned encoder alone
already provides a sufficiently structured state space for control, independent of
any learned dynamics.

\paragraph{Setup.}
We construct a control variant in which the agent never observes raw pixels.
Instead, each observation $o_t$ from the real environment is passed through
the frozen latent encoder $E_\xi$, and the agent receives only the resulting latent
representation
\[
z_t = E_\xi(o_t) \in \mathbb{R}^{d_z}.
\]
Crucially, all other components of the environment remain unchanged: the transition
function, reward function, and termination conditions are those of the true
environment. The world model predictor $f_\psi$ is not used in this setting.

The agent is the same action-chunked Flow Q-Learning (FQL) architecture introduced
in \S\ref{sec:fql-bg}, operating with a fixed chunk length $h=5$. Training proceeds
in two phases: offline pretraining on the encoded dataset, followed by interaction
with the real environment. This variant is a control: if the agent performs well
with latent observations and true dynamics, then later degradation under model
dynamics can be attributed to the predictor $f_\psi$ rather than to the encoder or
policy architecture. The quantitative result is reported in \S\ref{sec:exp-wm-as-env}.

% ---------------------------------------------------------------------
\section{Offline-to-Online Reinforcement Learning inside World Model}
\label{sec:wmep-exp1}

We now evaluate the most direct substitution: replacing the real environment with
the learned latent world model $\widetilde{\mathcal{E}}$ while keeping the rest of
the reinforcement learning pipeline unchanged.
If the world model provides a faithful approximation of the true environment
dynamics, then replacing the environment should yield a standard offline-to-online
improvement signal. The agent would effectively receive additional on-policy data,
allowing it to refine its policy beyond the offline dataset in the same way as it
would with real interaction.

\paragraph{Setup.}
Each episode begins by sampling an initial latent state from the empirical
initial-state distribution derived from the dataset. The agent then interacts with the world model for up to 40 world-model steps per episode. Each
step corresponds to one latent transition under the predictor $f_\psi$, and the
agent receives chunked transitions exactly as in the real environment.
Training begins from the offline checkpoint used throughout the corresponding
evaluation, ensuring that the policy is already well formed before online
interaction in the world model begins.

\paragraph{Failure mechanism.}
Although the predictor $f_\psi$ is locally accurate on held-out
one-step prediction, this does not imply stability under multi-step
rollouts. The model is trained only to predict the
next state, without any constraint enforcing long-horizon consistency with the data
distribution.
As a result, even small unbiased errors accumulate over long rollouts of up to
40 steps. This gradual drift moves the latent trajectory into regions that are
poorly represented in the training data. Once outside this region, both the critic
and actor are forced to extrapolate.
The critic, which is trained via temporal-difference updates, begins to evaluate
states that are effectively off-manifold. Because these states are not grounded in
real data, value estimates become increasingly unreliable. The actor then exploits
these inaccuracies by selecting actions that maximise the erroneous critic signal,
leading to policies that perform well inside the world model but fail when returned
to the real environment.
The corresponding real-environment evaluation is reported in
\S\ref{sec:exp-wm-as-env}. This comparison isolates the effect of replacing true
environment transitions with predicted ones while leaving the policy architecture
and latent representation fixed.

% ---------------------------------------------------------------------
\subsection{Latent States as Anchors}
\label{sec:wmep-exp2}

The naive WM-environment setup in \S\ref{sec:wmep-exp1} uses long
predicted rollouts, which makes compounding prediction error a likely
obstacle. Although the predictor can be accurate over a single step,
small errors accumulate as rollouts become longer, eventually producing
latent states that lie far from the distribution seen during
world-model training.
A natural mitigation is therefore to restrict how long the agent is
allowed to rollout purely under predicted dynamics. Rather than
executing long uninterrupted rollouts, we periodically attempt to
``re-anchor'' the trajectory to a latent state drawn from the offline
dataset.

\paragraph{Anchor mechanism.}

Let $\mathcal Z_{\mathrm{off}}$ denote the set of latent states present
in the offline dataset. Every $L_b=3$ world-model steps, the current
latent state is compared against this dataset and matched to its
nearest neighbour. If the neighbour is sufficiently similar under both
cosine similarity and Euclidean distance, the rollout is terminated and
restarted from that dataset state:

\begin{equation}
\Pi^{\epsilon}_{\mathcal Z_{\mathrm{off}}}(z)=
\begin{cases}
\hat z
&
\text{if }
\cos(z,\hat z)\ge \tau_{\cos}
\;\text{and}\;
\|z-\hat z\|_2 \le \tau_{L_2},
\\[4pt]
z
&
\text{otherwise},
\end{cases}
\label{eq:wmep-anchor-proj}
\end{equation}
where $\tau_{\cos}=0.95$ and $\tau_{L_2}=4.0$.

Intuitively, this mechanism attempts to keep trajectories close to
regions of latent space that are supported by real data. Rather than
allowing prediction errors to accumulate indefinitely, rollouts are
frequently reset to states that the world model has actually observed
during training.
Importantly, only the rollout initialisation is modified. The replay
buffer still stores the raw world-model transitions, meaning the agent
continues learning from predicted states rather than dataset states.
Anchoring successfully limits the most extreme forms of predictor
drift. However, the generated transitions are still produced by the
world model rather than the real environment, and therefore inherit any
errors present in the predictor.
We hypothesise that the primary issue is not the starting state of a
rollout but the transitions that are added to the replay buffer. Even
short rollouts can enter regions where the predictor is inaccurate,
producing transitions that do not correspond to valid environment
dynamics. Once these transitions are incorporated into the critic's
training distribution, value estimation becomes increasingly
unreliable.
This variant therefore tests whether limiting drift alone is sufficient when the
replay buffer still contains synthetic transitions. The sparse- and dense-reward
results are reported together with the other environment-substitution variants in
\S\ref{sec:exp-wm-as-env}.

% ---------------------------------------------------------------------

\subsection{Joint Training World Model and Policy}
\label{sec:wmep-exp3}

The previous experiments treat the world model as fixed. An
alternative is to adapt the predictor during reinforcement learning so
that it becomes more useful for control.
The base predictor is trained purely to predict
future latent states; however, a model that is accurate in a
prediction-error sense is not necessarily the model that best supports
value learning. If the world model is ultimately going to serve as the
agent's environment, it may be beneficial to optimise it directly for
the downstream RL objective. This follows the idea of value-aware
model learning, which trains dynamics models to minimise error in the
value function rather than raw prediction error
\cite{farahmand2017value}.

\paragraph{Joint objective.}

During offline training, the predictor $f_\psi$ is updated alongside
the policy using a combination of its original prediction loss and an
additional value-consistency objective:

\begin{align}
\mathcal L_{\mathrm{WM}}(\psi)
&=
\mathcal L_{\mathrm{pred}}(\psi)
+
\beta(t)\,
\mathcal L_{\mathrm{value}}(\psi),
\label{eq:wmep-joint}
\\
\mathcal L_{\mathrm{pred}}(\psi)
&=
\mathbb E\!\left[
\frac{\|f_\psi(z,a)-z'\|_2^2}
{\|z'\|_2^2+\varepsilon}
\right],
\label{eq:wmep-pred}
\\
\mathcal L_{\mathrm{value}}(\psi)
&=
\mathbb E\!\left[
\Big(
Q_{\bar\theta}(z,a)
-
\big(
r+\gamma m
Q_{\bar\theta}(f_\psi(z,a),a')
\big)
\Big)^2
\right],
\label{eq:wmep-value}
\end{align}
where

\[
a'
\sim
\pi_\theta\!\left(
\cdot
\mid
\mathrm{sg}[f_\psi(z,a)]
\right).
\]
The first term is the original prediction objective. The second
encourages the predictor to generate next states that are consistent
with the critic's Bellman targets.

Intuitively, if the critic believes a transition should have a
particular value, the predictor is rewarded for producing latent states
that preserve that value estimate. The weighting coefficient
$\beta(t)$ is gradually increased from $0$ to $0.1$ over the first
$50\,000$ optimisation steps so that value-learning signals do not
dominate before the critic has stabilised.

\paragraph{Diagnostic component swap.}

Any change in performance could in principle originate from either component: the
jointly trained policy, the jointly trained predictor, or their
interaction. To separate these possibilities, we evaluate a
component-swap diagnostic in Chapter~\ref{chap:experiments}. One
configuration pairs the original offline policy with the jointly trained
predictor; the other pairs the jointly trained policy with the original
frozen predictor. The numerical results are reported in
\S\ref{sec:eval-joint}, since this diagnostic is part of the evaluation
rather than a separate training method.

A plausible explanation is that the value-consistency objective
introduces a form of model exploitation.
The intended solution of Eq.~\eqref{eq:wmep-value} is for the predictor
to generate latent states that are both accurate and value-consistent.
However, the loss also admits undesirable solutions. For example, the
predictor can reduce the Bellman residual by mapping many inputs to a
small set of critic-preferred latent states, even when those states do
not correspond to the true dynamics.
In this failure mode, the value objective becomes much easier to optimise than
the prediction objective. Rather than learning increasingly accurate
dynamics, the predictor may gradually distort its outputs toward
regions that appear favourable to the critic.
This mechanism cannot be observed directly from success rates alone, which is why
the component-swap diagnostic is included in the evaluation. It is also consistent
with prior work such as MOPO \cite{yu2020mopomodelbasedofflinepolicy}, which keeps
the world model fixed and instead accounts for model uncertainty during value
learning.

\subsection{Policy Optimisation with Uncertainty-Penalised Rewards}
\label{sec:wmep-exp45}

The previous experiments suggest that simply replacing the real
environment with the world model is insufficient. A common response in
model-based offline RL is to retain the learned dynamics model but
penalise transitions that are likely to be unreliable. Methods such as
MOPO and COMBO follow this principle by reducing the reward assigned to
states where the model is uncertain
\cite{yu2020mopomodelbasedofflinepolicy,yu2021combo}.

To test whether such penalties can stabilise WM-environment training,
we augment the anchored dense-reward setup of
\S\ref{sec:wmep-exp2} with the MOPO uncertainty penalty:

\begin{equation}
    \tilde r'_t \;=\; \tilde r_t \;-\; \lambda\, u(z_{t+1}, a_t),
    \label{eq:wmep-mopo-reward}
\end{equation}
where $u(\cdot)$ estimates the uncertainty of the predicted next state.

\paragraph{Uncertainty estimators.}

We consider two commonly used proxies for uncertainty.

\textbf{Ensemble disagreement.}
Three independently initialised predictors are trained in parallel. If
their predictions disagree, the transition is treated as uncertain:

\begin{equation}
    u_{\text{ens}}(z, a) \;=\; \Big\|\,\sigma_k\!\big[f_\psi^{(k)}(z, a)\big]\,\Big\|_2,
    \label{eq:wmep-mopo-ens}
\end{equation}
This is the standard ensemble-based uncertainty estimate used in many
model-based RL methods.

\medskip

\textbf{Nearest-neighbour distance.}
Instead of maintaining multiple models, uncertainty can also be
estimated from how far a predicted latent lies from the offline
dataset:

\begin{equation}
    u_{\text{nn}}(z') \;=\; 1 \;-\; \max_{z\in\mathcal{Z}_{\text{off}}}\;
        \tfrac{\langle z, z'\rangle}{\|z\|\,\|z'\|},
    \label{eq:wmep-mopo-nn}
\end{equation}
Predictions that are dissimilar to every offline latent receive a
larger penalty.

The two penalties test whether uncertainty estimates can separate useful imagined
transitions from harmful ones. One possible failure mode is that neither quantity is measuring exactly
what the agent needs. The nearest-neighbour penalty measures distance
from the offline dataset rather than prediction error. In offline
goal-reaching datasets, successful trajectories may naturally move
toward regions of latent space that are sparsely represented in the
behaviour data. Such states are not necessarily unreliable, yet they
receive the largest penalties. The agent may therefore be discouraged
from taking actions that genuinely advance the task.
The ensemble estimate avoids this particular issue, but disagreement between
predictors may still fail to identify the specific errors that corrupt value
learning. The numerical comparison is reported in \S\ref{sec:exp-wm-as-env}.

\paragraph{Relationship to the next chapter.}

The next chapter uses a different source of grounding. The offline
critic is trained entirely on real transitions and therefore contains
information about which regions of latent space are supported by the
dataset. In contrast, both uncertainty estimators attempt to infer
support indirectly through either model disagreement or dataset
similarity.
Rather than
explicitly penalising uncertainty, the world model is queried only for
short-horizon predictions and candidate trajectories are ranked directly
by the offline-trained value function. This keeps the learned dynamics in a
proposal role while leaving the final judgement to a critic trained on real data.

% ---------------------------------------------------------------------

\section{Design Risks of Optimisation Inside World Model}
\label{sec:wmep-synthesis}

The variants in this chapter are designed around three risks that arise when a
learned dynamics model is used as a direct replacement for the environment:

\begin{enumerate}
\item \textbf{Prediction errors accumulate over long rollouts.}

The predictor is trained using a one-step objective, yet online RL
requires trajectories spanning tens of world-model steps.
Across such horizons, small inaccuracies compound and gradually
move trajectories away from regions where the predictor is known to
be reliable.

\item \textbf{Learning from synthetic transitions can corrupt value estimates.}

Regardless of whether anchoring or uncertainty penalties are used,
the critic is ultimately trained on transitions generated by the
world model. If those transitions are systematically inaccurate,
the resulting value function may assign incorrect values to
real-environment behaviours.

\item \textbf{Updating the world model using RL objectives is risky.}

Value-aware objectives can in principle distort the predictor away from faithfully
modelling the environment. The component-swap diagnostic in
\S\ref{sec:eval-joint} is included to test whether such distortion comes from the
predictor itself or from the policy trained alongside it.

\end{enumerate}

These risks motivate the design choices of the following chapter:
restricting the world model to short-horizon imagination, keeping the critic
grounded in real data, and avoiding RL-driven updates to the predictor.

% ---------------------------------------------------------------------

\section{Summary}
\label{sec:wmep-summary}

This chapter defined the direct environment-substitution variants and the risks
they are intended to test. The quantitative comparison is collected in
\S\ref{sec:exp-wm-as-env} (Table~\ref{tab:owm-summary}). The remaining chapters
therefore treat the world model not as an environment to learn inside, but as a
short-horizon source of candidate futures whose value is judged by models trained
on real data.
